\chapter{Some notes from Humphreys}
\section{Pre - varieties and varieties}
\subsection{Projective varieties}
\begin{dang}
    \begin{definition}
        Let $V$ be a vector space, we define its associated projective variety by 
        \begin{equation}
            \bbP(V) = (V\setminus0 )/ \sim 
        \end{equation}
        where we let $v\sim w\iff v=\lambda w$ for some $\lambda\in k^\times$. Suppose that $\dim(V)=n+1$, 
        then by choosing a basis of $V$, we identify $V=k^{n+1}$ and we write $\bbP(V)=\bbP^{n}$.
    \end{definition}
\end{dang}
 
We define the Zariski topology on $\bbP$ by letting a closed set to be a locus of a homogeneous polynomial. So 
$X\subset \bbP^n$ is closed if and only if $X=V(S)$ for some collection $S\subset k[X_0,\dots,X_n]^h$ of homogeneous polynomials. 
We have a set of isomorphisms ($i=0,\dots,n$):
\begin{equation}
  \phi_i:  k[X_1,\dots,X_n]\rightarrow k[X_0,\dots, X_n]^h: g\mapsto
   X_i^{\deg(g)}g(\frac{X_0}{X_i},\dots,\frac{X_{i-1}}{X_i},\frac{X_{i+1}}{X_i}\dots,\frac{X_n}{X_i}).
\end{equation}
The inverse of this map takes $\phi_i:h(X_0,\dots, X_n)\mapsto h|_{X_i=1}$. In fact 
if $h=\sum_{\alpha\in\bbZ_{\geq0}^{n+1}}a_\alpha X^\alpha\in k[X_0,\dots, X_n]$ is homogeneous of degree $d$, then
\begin{equation}
    \begin{split}
        \psi_i\phi_i(h)&=X_i^dh(\frac{X_0}{X_i},\dots,\frac{X_{i-1}}{X_i},1,\frac{X_{i+1}}{X_i},\dots,\frac{X_n}{X_i})\\
        &=\sum_\alpha a_\alpha X_1^{\alpha_1}\cdots X_{i-1}^{\alpha_{i-1}} X_i^{d-\sum_{j\neq i}\alpha_j}
        X_{i+1}^{\alpha_{i+1}}\cdots X_n^{\alpha_n}\\
        &=h,\;(\text{because }\sum_j\alpha_j=d).
    \end{split}
\end{equation} 
Conversely 
\begin{equation}
        \phi_i\psi_i(g) = X_i^dg(\frac{X_0}{X_i},\dots,\frac{X_n}{X_i})\bigg|_{X_i=1}=g.
\end{equation}
A homogeneous coordinate of $\bbP^n$ is a map 
\begin{equation}
    f_i: U_i=\{[x_0,\dots,x_n]:x_i\neq0\}\rightarrow k^n:[x_0,\dots, x_n]\mapsto 
    (\frac{x_0}{x_i},\dots, \frac{x_{i-1}}{x_i},\frac{x_{i+1}}{x_i},\dots, \frac{x_n}{x_i}).
\end{equation}
Since $U_i = \bbP^n\setminus V(x_i)$, it's an open set by definition. Furthermore, the map $f_i$ is an 
isomorphism. In fact, let $Y\subset U_i$ be a closed set, then $\overline{Y}$ is closed in $\bbP^n$, so there exist 
$S\subset k[X_0,\dots, X_n]^h$ such that $\overline{Y}=V(S)$. So we claim that $f_i(Y)=V(\psi_i(S))$. In fact, we have 
\begin{equation}
    \begin{split}
        x=(x_0,\dots,x_{i-1},x_{i+1},\dots,x_n)\in V(\psi_i(S)) &\iff \phi_i(h)(x) = 0\;\forall h\in S\\
        &\iff h[x_1,\dots,x_{i-1},1,x_{i+1},\dots, x_n]=0\;\forall h\in S\\
        &\iff hf_i^{-1}(x)=0\;\forall h\in S\\
        &\iff f_i^{-1}(x)\in V(S)\cap U_i = Y\\
        &\iff x\in f_i(Y).
    \end{split}
\end{equation}
Conversely, let $Z=V(T)\subset k^n$. We want to show that $f_i^{-1}(Z)\subset U_i$ is closed. We claim 
that $f_i^{-1}(Z)=V(\phi_i(T))\cap U_i$, which can be proved in the sames lines as before.\par 

\begin{dang}
    \begin{lemma}
        (Affine characterization for closedness). A subset $X\subset\bbP^n$ is closed if and only if 
        $X\cap U_i$ is closed in $U_i$ for all $i=0,\dots,n$.
    \end{lemma}
\end{dang}

\begin{dang}
    \begin{lemma}
        Consider the map 
        \begin{equation}
            \begin{split}
                 \varphi:&\bbP^{m}\times \bbP^n\longrightarrow\bbP^{(m+1)(n+1)-1}\\
                 &([x_0,\dots, x_m],[y_0,\dots,y_n])\longmapsto [\dots,x_iy_j,\dots]
            \end{split}
        \end{equation}
        where we use $[\dots,x_iy_j,\dots]$ to abreviate $[x_0y_0,\dots, x_0y_n,\dots, x_my_0,\dots, x_my_n]$. 
        We claim that $\varphi(\bbP^m\times\bbP^n)\subset\bbP^{(m+1)(n+1)-1}$ is a closed embedding, so it's 
        an isomorphism onto its image.
    \end{lemma}
\end{dang}
Proof: Use $U_i,V_j, W_{kl}$ to denote the affine open charts of 
$\bbP^{m},\bbP^n$ and $\bbP^{(m+1)(n+1)-1}$ respectively. The strategy is to show that 
$\varphi(U_i\times V_j)$ is closed in $W_{ij}$. However, if $(x,y)\in U_i\times V_j$, then 
$\varphi(x,y)=[\dots,x_ky_l,\dots]$, where we assume $x_i=1,y_j=1$. Now, we observe that the image is of the form 
$[\dots,u_{kl},\dots]$, where $u_{kl}=x_ky_l$. Then we conclude that $\varphi(U_i\times V_j)$ is the locus 
\begin{equation}
   f_{ij}^{-1}( \{[u_{kl}:(k,l)\neq(i,j),\;u_{kl}-u_{kj}u_{il}=0,\;\forall k,l]\}) \overset{\text{closed}}{\subset}W_{ij}.
\end{equation}
It remains to show that 
$\varphi$ is injective. To do this, it's sufficient to see that $\varphi|_{U_i\times V_j}$ 
(viewed as a map to $\varphi(U_i\times V_j)$) can be inverted. The equation is just 
\begin{equation}
    [\dots, u_{kl},\dots]\mapsto ([u_{0j},\dots, u_{mj}],[u_{i0},\dots, u_{in}]).
\end{equation}
\qed
\begin{dang}
    \begin{corollary}
        If $X\subset \bbP^m, Y\subset\bbP^n$ are closed subsets, then $\varphi(X\times Y)\subset\bbP^{(m+1)(n+1)-1}$
        is a closed subset.
    \end{corollary}
\end{dang}

\subsection{Varieties}
